\section{A classical approximate-query baseline}
\label{sec:sketch-baseline}
The exact dimension theorem does not imply that approximate residual queries
need fourth-order storage. Here is an explicit transfer from an existing
matrix-sketch guarantee. This is a corollary, not a new sketching algorithm.

Let $\operatorname{svec}:\operatorname{Sym}_d\to\mathbb R^p$ be the
Frobenius isometry which stores diagonal entries unchanged and off-diagonal
entries multiplied by $\sqrt2$. Form $A\in\mathbb R^{n\times p}$ with
row $i$ equal to $\operatorname{svec}(H(x_i))^\top$. Store $S_y$ and
$\overline H$ exactly, and let $C\in\mathbb R^{\ell\times p}$ satisfy
\[
 \|A^\top A-C^\top C\|_2\le\eta.
\]
Define
\[
 \widehat G(B,c)=S_y-
 \operatorname{svec}^{-1}\!\left(\frac{C^\top C}{2n}
                       \operatorname{svec}(B)\right)
 -\frac c2\overline H.
\]
\begin{proposition}[Uniform deterministic query transfer]
For every $R\ge0$, simultaneously for every $\|B\|_F\le R$ and
$c\in\mathbb R$,
\[
 \|G(B,c)-\widehat G(B,c)\|_F\le\frac{R\eta}{2n}.
\]
The guarantee therefore remains true for queries chosen adaptively from the
summary or the entire data, without an independence assumption.
\end{proposition}
\begin{proof}
The isometry and the empirical query formula give
\[
 \operatorname{svec}(G-\widehat G)
 =-\frac{A^\top A-C^\top C}{2n}\operatorname{svec}(B).
\]
The operator norm inequality proves the result, pointwise and uniformly.
\end{proof}

Frequent Directions gives, for every integer $0\le k<\ell\le p$,
$\eta\le\|A-A_k\|_F^2/(\ell-k)$, where $A_k$ is a best rank-$k$
approximation \cite[Theorem 1.1]{p7ghashami2015}. Thus the query error is
at most $R\|A-A_k\|_F^2/[2n(\ell-k)]$ using $\ell p+2p$ stored real
coordinates for $C,S_y,\overline H$, plus a count and implementation
workspace. This count excludes floating-point rounding and is not a
claim of optimality for the polynomially constrained feature rows.
The cited theorem concerns arbitrary matrices; any approximate lower
bound from that paper cannot automatically be restricted to rows
$\operatorname{svec}(xx^\top-I)$.

This baseline has no distribution assumption, but its tail bound can be
large. It neither bounds out-of-sample risk nor proves that the selected
growth direction is useful. A next research target is a specified
memory--accuracy tradeoff on these polynomial feature rows, with a
separate population-error and optimization analysis. No sketching
implementation or new training experiment was run in Phase 7.
